%% kapa: exam=1 round=26 year=2015 period=1 form=A subject=economics subject_ko=경제원론 seq=22 number=62
%% source: Q-Net 감정평가사 기출문제 (cst00309) artlSeq=5138001 file=extracted/1008/5138001/제26회_감정평가사 1차_기출/제26회_감정평가사 1차_기출/1교시_감평_A형.pdf
%% conversion: kordoc 4.24.51 -> tools/gen1.py (auto); review=proofread-sonnet; flags=none
\begin{question}{62}
쌀과 컴퓨터만 생산하는 국가의 생산량과 가격이 다음과 같다. 2013년을 기준연도로 할 때 2014년의 실질GDP와 실질GDP성장률은?
\par{\small\centering\begin{tabular}{|c|c|c|c|c|}\hline
 & \multicolumn{2}{c|}{쌀} & \multicolumn{2}{c|}{컴퓨터} \\ \cline{2-5}
 & \makecell{가격\\(원)} & \makecell{생산량\\(가마)} & \makecell{가격\\(원)} & \makecell{생산량\\(대)} \\ \hline
2013년 & 10 & 50 & 30 & 100 \\ \hline
2014년 & 15 & 100 & 50 & 200 \\ \hline
\end{tabular}\par}
\begin{choices}
  \item 3,500원, 100 \%
  \item 3,500원, 228.6 \%
  \item 7,000원, 100 \%
  \item 7,000원, 228.6 \%
  \item 11,500원, 64.3 \%
\end{choices}
\end{question}
\answer{3}
